\chapter{Conclusion}

\label{chapter_conclusion}

This thesis has developed and presented a new polynomial expansion technique, the analytically modified integral transform expansion (AMITE), designed to facilitate the detailed and exact analysis of a broad class of nonlinear functions.  The main motivation for this development was the need to formulate an improved polynomial expansion technique suited for application to the nonlinear activation functions often employed in neural networks.  To frame the gaps addressed in current research, a number of existing polynomial expansions were reviewed along with their applications to neural networks.  The studies that were reviewed focused on
addressing well-known difficulties in the verification, explainability, and
security of neural networks.  These challenges were further
motivated by listing just a few neural network applications of interest,
e.g., applications in health-care and defense, which have driven
the artificial intelligence
engineering community to treat these learning machines with analytical rigor.

While the existing techniques for polynomial expansion were found to provide useful properties individually, many of these properties were mutually exclusive even when using the state of the art techniques.  For example, expansions which provide exact formulas, such as the classical Taylor series, do not also provide an adjustable region of convergence.  Expansions which provide an adjustable region of convergence, e.g., numerical methods, do not provide exact equations for the errors.  With these gaps established and a set of target functions for expansion identified, a new technique was developed to provide all the desired properties of the many previous expansions considered, including those properties which were otherwise mutually exclusive.  In Chapter \ref{chapter_introduction}, the desired properties were reviewed and a summary of available expansions and their properties was presented.   Expansions that were considered include Taylor, Fourier-Legendre, Chebyshev, asymptotics, and several numerical techniques. The desired properties that were considered are given as follows.
\begin{enumerate}
	\itemsep-0.5em 
	\item Exact coefficient formulas
	\item Explicit exact error formulas
	\item Monic form
	\item Adjustable domain
	\item Handling of undefined derivatives
	\item Consistency of method
\end{enumerate}

The main research contribution of this thesis was therefore to develop a polynomial expansion which has all the advantageous properties listed above and which can be applied to nonlinear activation functions of neural networks.  In addition to this main motivation, significant auxiliary motivations were also established.  The nonlinear functions which are used in neural network activation units also arise naturally in the modeling and the study of other nonlinear systems.  The analysis of neural network nonlinear activation functions therefore stands to benefit many other domains as well.  Two diverse examples considered were shallow water traveling waves and nonlinear circuits, both of which involve similar nonlinear behavior to the activation functions of neural networks.  For the benefit of all of these applications, the explicit expansions of four functions, $\sech(v)$, $\tanh(v)$, $\relu(v)$, $\sech^2(v)$, were developed and analyzed.  Exact error formulas were derived and presented for the first time, along with novel and useful closed form error approximations.

The main in-depth case study selected to demonstrate the utility of the expansion was performed on a Multi-Layered Perceptron (MLP) neural network.  In this case study, two sets of experiments were conducted.  First, a series of fundamental experiments, denoted experiment group A, were performed on a small, two input, three hidden node, single output neural network.  Second, a series of further experiments, denoted experiment group B, were performed on a population of neural networks of varying size, up to 125 hidden neurons.  The purpose of experiment group A was to perform fundamental analysis on a network small enough to be analyzed in significant detail.  The purpose of experiment group B was to assess the scalability of the technique for practically sized neural networks which would be employed in neural control applications.  To achieve this, experiment group B also required optimized programming via the use of graphic processing units and other optimization techniques including custom optimized memory management for the implementation of some computationally expensive mathematical functions.

Of experiment group A, the experiments executed and the conclusions drawn therefrom are as follows.  The effect of the entropy of the stimulus signal on the accuracy of the conformance test was analyzed and it was found that a high entropy stimulus is required to achieve a reliable test.  This result agrees with known results from the characterization of nonlinear circuits.  Analysis was performed to extract the input-output manifolds of three pairs of networks including symmetrically conformant, approximately conformant, and non-conformant networks.  It was found that symmetrically conformant networks yield zero error between their AMITE coefficients, approximately conformant networks have small errors between their AMITE coefficients, and in contrast, non-conformant networks have large errors between their AMITE coefficients.  These results demonstrate that the AMITE coefficients represent a space wherein pairs of networks can be effectively conformance tested with symmetrically conformant networks being treated as conformant, as is desired.

To determine how reliable such a conformance test is, a receiver operating characteristic (ROC) curve was computed for varying signal to noise ratios which the measurements of the network output were subject to.  The ROC curve was employed to identify an optimal error threshold, $\varepsilon = 1$, to pair with the inverse hyperbolic sine scaled error metric used in experiment group A.  The optimal error threshold was then used to analyze test accuracy over a population of 50 defective and 50 conformant networks with respect to measurement SNR.  It is noted that the networks in this experiment group were treated as black box systems, with no initial guess for the expected weights of the network under test provided to the test algorithm.  Even with these challenging constraints, we find the conformance test reliable for $> 80\ \mathrm{dB}$ measurement SNR.

In experiment group B, a more realistic and practical set of tests were conducted.  A population of 1890 neural networks with hidden units ranging from 5 to 125 and inputs ranging from 2 to 4 was instantiated.  In this test, the fact that the expected weights for the conformance test are known was exploited and these values were used as a starting point to search for the weights of the replicated network.  Further, while experiment A used inverse hyperbolic sine error scaling, experiment B used mean logarithmic error scaling.  This change was made because the larger size of the network necessitated a more sensitive error metric since changes in individual weights among many weights cause more subtle changes in the output manifold than changes in individual weights among few weights.  Conformance test accuracy was analyzed with respect to measurement SNR and it was found that for experiment group B, a $90\%$ accurate test could be achieved even for measurement SNRs approaching $0\ \mathrm{dB}$.

To assess scalability, runtimes were tabulated with respect to the number of hidden units and the number of inputs.  It was found that the conformance test procedure terminated in under 13 seconds for networks having up to 125 hidden neurons and 4 hidden nodes.  Further, as expected from the structure of the algorithm for generating the AMITE coefficients of the MLP as presented in Chapter \ref{chapter_improved_expansions}, the time to run the conformance test was found to increase linearly with respect to the number of hidden units and factorially with respect to the number of inputs.  This implies that the conformance test procedure is best suited for networks with just a few inputs but can be scaled to many hidden nodes while still providing practical run times.  The results of experiment group A and experiment group B therefore demonstrate that the AMITE coefficients provide a polynomial representation by which MLPs can be reliably conformance tested even when their true outputs are obscured by measurement noise and the values of their weights are inaccessible.

With the utility of the AMITE polynomials demonstrated for neural networks in the MLP case study, their general utility was then demonstrated in two complementary case studies in traveling shallow water waves and nonlinear circuits.  First, a basic problem in traveling shallow water waves was treated to illustrate use of the AMITE technique to break up an integral in a manner analogous to the classical application of Taylor expansion for the same purpose.  An expression for the mean height of a wave at some point in space and time given some uncertainty about its measured wavenumber was derived and found to be analytically intractable and resistant to application of the Taylor series, therefore exhibiting the same difficulties which are encountered when attempting to treat neural network nonlinearities with existing expansions.  Using the AMITE polynomials, the integrand was expanded successfully such that the problem could then be solved analytically.  Agreement was shown between analytical and numerical results, indicating that both could be trusted and that the system had been appropriately characterized.

In a deeper case study, two nonlinear circuits, a nonlinear amplifier without D.C. bias and a nonlinear amplifier with D.C. bias, were characterized.  The odd harmonics of the amplifier without bias were expressed as AMITE polynomials adjusted by a set of integrals which were then solved exactly.  The exact values of the integrals for each polynomial coefficient, for each harmonic at output, were found to form a family of infinitely many novel integer sequences, none of which have never been cataloged before in the Online Encyclopedia of Integer Sequences (OEIS). For the nonlinear amplifier with input bias, its bias at output was computed as a function of the bias and gain of its input signal, accounting exactly for the mixing effects between the contributions of the input gain and bias to the output bias.  The results were again expressed via AMITE polynomials adjusted by a set of integrals which were then solved exactly.  The exact values of these integrals form infinitely many additional families of novel integer sequences, each family therein containing infinitely many sequences which have never before been cataloged in the OEIS.  This case study therefore served to further establish that the AMITE polynomials are a useful new approach for exactly characterizing the behavior of nonlinear circuits.

The AMITE polynomials and the results in the case studies for neural networks, shallow water waves, and nonlinear circuits have therefore opened many directions for new work in all these fields and more.  These directions for future work were reviewed in Chapter \ref{chapter_future_work} and a technical approach for extending the case studies to even more applications than were presented here has been outlined therein.  The technical approach rests on three new topics for fundamental research which would enable solving an noteworthy array of fundamental and practical problems in the analysis of nonlinear systems.  These enabling fundamental research topics are (1) iterative / recursive application of the AMITE polynomials, (2) methods of accounting for nonlinear expansion error propagation through several layers of nonlinearity, and (3) methods for extracting more general analytic forms, such as Meijer-G or Fox-H forms from the AMITE polynomials and from the results of problems solved analytically via the AMITE polynomials.  These fundamental mathematical advancements, building off the advancements herein, would enable impactful new directions for research in neural network verification, neural network interpretability and explainability, analysis of nonlinear waves, and nonlinear circuit design and analysis, including applications in radar systems engineering, music and audio signal processing, and communications.

From all the case studies, results thereof, and directions for new work thereby opened, it is established that the AMITE polynomials represent a powerful new dimension of expansion methods suitable for detailed and exact analysis of a wide class of nonlinear systems.  In closing, it is noted that the AMITE polynomials and the case studies wherein they are applied provide some deep insight into the commonalities underlying many diverse types of nonlinear systems and behaviors.  The systems that the AMITE polynomials can characterize are diverse and have many different underlying processing mechanisms, with the few examples considered herein spanning artificial models of neural processing (i.e., neural networks), ocean wave dynamics, and distortions induced in a wide variety of electrical circuits, themselves an effect of fundamental physics constraints which govern the behavior of those circuits.  The fact that such a wide variety of phenomena can be characterized by a common method which exhibits striking parallels in its application to the unique peculiarities of each shines light on the common motifs which underlie the behavior of nonlinear systems in general, and serves as a small example of the power of mathematical abstraction to describe these common motifs in a manner which allows their employment as reusable tools in the solution of practical engineering problems.

